\subseccion{Conferencista Internacional}

%\vspace{-0.5cm}
%\subtitulos{Conferencista Internacional}
%===============================================================================================
%								6ta conferencia (internacionales)
%===============================================================================================

\bandera{ES}
\confname{Dra. Cristina González Montelongo}
\confsch{Universidad de La Laguna}
\conffac{}
\confjob{Profesora e investigadora\, \href{https://orcid.org/0000-0001-9236-890X}{\aiOrcid}\, \href{https://www.facebook.com/paulino.aguilar.7568596/videos/1092630373338066/}{\faTv}}
\confimgone{images/speakers/Dra. Cristina González Montelongo}
\confimgtwo{images/speakers/logo ull}
\conftext{4}{3}{1}{%
	La Dra. Cristina González Montelongo es bióloga con Máster Universitario de Biodiversidad Terrestre y Conservación en Islas, y Doctora por la Universidad de La Laguna. En 2012 se inició en la investigación botánica a través de una Beca de Colaboración, que dio paso a dos Proyectos de Innovación Educativa.  Ha alternado sus estudios de máster con el trabajo para empresas del sector medioambiental. Permitiéndole desarrollar actividades de educación ambiental, informes de impacto ambiental, Seguimientos de Especies Amenazadas (SEGA), actualización de las bases de datos de especies silvestres y de especies amenazadas del Gobierno de Canarias, entre otras actividades. Su tesis doctoral se centró en el estudio del monteverde canario empleando los líquenes epífitos como organismo modelo. Estos estudios los alternó con el trabajo como técnico en el Herbario TFC (Herbario Institucional de la ULL, SEGAI), encargándose desde entonces del mantenimiento del mismo y su informatización, además de facilitar mejoras en la gestión de las colecciones y en las propias instalaciones, sin abandonar en ningún momento las actividades de divulgación, docencia (reglada y no reglada), investigación y asesoramiento técnico a investigadores y empresas del sector medioambiental, a través de su trabajo en el Herbario TFC. Su trabajo e investigaciones han resultado en 8 informes científico-técnicos, 14 artículos científicos, 1 artículo de divulgación científica, más de 30 contribuciones a congresos nacionales e internacionales y más de 30 participaciones en ferias, festivales y jornadas de divulgación, además de haber organizado un congreso científico y dos jornadas técnicas. Es Personal Técnico de Apoyo en el Servicio Herbario TFC del SEGAI, y miembro del grupo de investigación “Liquenología” del Departamento de Botánica, Ecología y Fisiología Vegetal. Actualmente, colabora con numerosos equipos de investigación y proyectos tanto de la ULL como nacionales e internacionales.%
}

\conference

\newpage
\begin{refsection}
\pagenblanc{no}
\begin{wrap}
	\chapter{Francisco Véliz}
	% If you're using TexStudio, this line helps you to identify the keynote/document/course in the editor's TOC
\end{wrap}
\cero
\abbrevnombs{Véliz-Romero}{Francisco}{Francisco-Gerardo}{Gerardo}
\tocauthor{Véliz-Romero Francisco Gerardo}% This function adds the name in the TOC
\author{%
	Véliz-Romero Francisco Francisco-Gerardo Gerardo
}
\ijaffiliation{%
	Asociación Mexicana de Productores de Carne AMEG, A. C.
}
% In the first {} you can add any text, and the second {} the email
% The function adds automatically an space between both items
\corresponding{Este es un ejemplo:}{exaple@gmail.com}
\title{Keynote's title}

\maketitle

\vspace{.5cm}
\begin{resumen}
	\blindtext
	\blindtext
	\blindtext
\end{resumen}

\keywords{i'll say, maximum four, keynotes, but, it doesn't, have, limits}

% I personally dont need the subsection's title in TOC, so i used the *
\subsection*{The subsection: Of the keynote}

\blindtext[1] \textcite{example}. \blindtext[1]

\blindtext[1] \parencite{example-5}

\subsection*{Other subsection}

\blindtext[1] \parencite{example-2}.

\blindtext[1] \textcite{example-1,example-2} \blindtext[1] \parencite{example-6}.

\subsection*{More sections?}

\blindtext[1]

\begin{itemize}
	\item[a.] \blindtext[1]
	\item[b.] \blindtext[1]
	\item[c.] \blindtext[1]
	\item[d.] \blindtext[1]
\end{itemize}

Altough i used itemize, it sould also support ennumerate environment too, though i havent checked it.

% Adding references that are not printed in the text don't break the document ;)
\nocite{example-3,example-4}

% if you dont want the bibliography title appearing in the toc, you can use "misubseccion"
% If you want the bibliography title appearing in the toc, change it to "subbibintoc"
\printbibliography[heading=misubseccion]
\end{refsection}

\newpage
\pagenblanc{alone}
\subseccion{Conferencista especial}

%\vspace{-0.5cm}
%\subtitulos{Conferencista Especial}

%===============================================================================================
%								Conferencia Especial
%===============================================================================================

\bandera{}
\confname{MVZ. Mario Guevara Acosta}
\confsch{Dirección General de Salud Animal}
\conffac{Comisión México–Estados Unidos para la Prevención de la Fiebre Aftosa y otras Enfermedades Exóticas de los Animales (CPA)}
\confjob{Coordinador Regional de la CPA Región Laguna II\, \href{https://www.facebook.com/share/v/1CghyqfprV/}{\faTv}}
\confimgone{images/speakers/MVZ. Mario Guevara Acosta}
\confimgtwo{images/speakers/logo cpa}
\conftext{4}{3}{1}{%
	\textbf{Generalidades del Gusano Barrenador}\\
	\textit{11 de septiembre de 2026}\\
	El MVZ. Mario Guevara Acosta es egresado de la Universidad Autónoma Agraria Antonio Narro, y ha ostentado el cargo de Profesional Ejecutivo de Servicios Especializados en la Dirección General de Salud Animal del CPA, bajo el Servicio Nacional de Sanidad, Inocuidad y Calidad Agroalimentaria (SENASICA) durante los años 2011 y 2012. En la actualidad, ejerce como Coordinador Regional de la CPA Región Laguna II desde el año 2025, y es presidente del Grupo Estatal de Emergencia en Sanidad Animal (GEESA), conformado por el Gobierno del Estado de Chihuahua. Impartiendo capacitaciónes sobre las medidas de prevención, vigilancia y control del Gusano Barrenador del ganado, así como de los procesos y protocolos que se deben seguir ante una contingencia sanitaria. Como parte de su trabajo como Coordinador Regional de la CPA Región Laguna II, ha realizado conferencias en universidades como la Universidad Juárez del Estado de Durango, la Universidad Autónoma Agraria Antonio Narro y la Unidad Regional Universitaria de Zonas Áridas de la Universidad Autónoma Chapingo, así como  curso-talleres y capacitaciones en la Unión Ganadera Regional de Durango, la Unión Ganadera Regional de Coahuila y la Asociación de Médicos Veterinarios Zootecnistas Especialistas en Bovinos de la Laguna; presentando información sobre el panorama actual del gusano barrenador, así como los riesgos, impactos y estrategias de control y erradicación establecidas para combatir esta problemática sanitaria.%
}

\conference
